\chapter{Fonctionnalités réalisées}

\section{Besoins fonctionnels}

Tous les besoins sont réalisés, y compris les besoins P3 (BF04, BF16, BF23).

{\small
\begin{xltabular}{\linewidth}{P{1.1cm}P{3.6cm}ccL}
  \caption{Besoins fonctionnels et leur réalisation}\label{tab:bf}\\
  \ligneentete
  \entete{Code} & \entete{Besoin} & \entete{Prio.} & \entete{État} & \entete{Réalisation} \\
  \endfirsthead
  \ligneentete
  \entete{Code} & \entete{Besoin} & \entete{Prio.} & \entete{État} & \entete{Réalisation} \\
  \endhead
  \bottomrule
  \endfoot
  BF01 & S'inscrire & P1 & Réalisé & \code{InscriptionServlet}, \code{AuthService.inscrire} \\
  BF02 & Se connecter & P1 & Réalisé & \code{ConnexionServlet}, \code{AuthService.authentifier} \\
  BF03 & Se déconnecter & P1 & Réalisé & \code{DeconnexionServlet} (POST, invalidation de la session) \\
  BF04 & Consulter et modifier son profil & P3 & Réalisé & \code{ProfilServlet}, \code{AuthService.modifierProfil}, \code{changerMotDePasse} \\
  BF05 & Liste des produits disponibles & P1 & Réalisé & \code{CatalogueServlet}, \code{CatalogueService.listerProduits} (paginée) \\
  BF06 & Filtrer par catégorie & P1 & Réalisé & Paramètre \code{categorie} du catalogue \\
  BF07 & Rechercher par mot-clé & P2 & Réalisé & Paramètre \code{motCle} (nom et description, insensible à la casse) \\
  BF08 & Fiche détaillée d'un produit & P1 & Réalisé & \code{ProduitServlet}, \code{CatalogueService.trouverProduitActif} \\
  BF09 & Ajouter au panier & P1 & Réalisé & \code{PanierServlet}, \code{PanierService.ajouter} \\
  BF10 & Modifier la quantité & P1 & Réalisé & \code{PanierService.modifierQuantite} \\
  BF11 & Retirer un produit, vider le panier & P1 & Réalisé & \code{PanierService.retirer}, \code{vider} \\
  BF12 & Consulter le panier et son total & P1 & Réalisé & \code{client/panier.jsp}, \code{Panier.getTotal} \\
  BF13 & Passer une commande & P1 & Réalisé & \code{CommandeServlet}, \code{CommandeService.passerCommande} \\
  BF14 & Historique des commandes & P1 & Réalisé & \code{CommandeService.listerCommandesClient} (paginé) \\
  BF15 & Détail d'une commande & P2 & Réalisé & \code{CommandeService.trouverCommandeClient} \\
  BF16 & Annuler une commande « En attente » & P3 & Réalisé & \code{CommandeService.annulerCommandeClient} \\
  BF17 & CRUD des catégories & P1 & Réalisé & \code{AdminCategorieServlet}, \code{CatalogueService} \\
  BF18 & CRUD des produits & P1 & Réalisé & \code{AdminProduitServlet}, \code{CatalogueService} \\
  BF19 & Gérer le stock & P1 & Réalisé & \code{CatalogueService.modifierStock} \\
  BF20 & Consulter et filtrer les commandes & P1 & Réalisé & \code{AdminCommandeServlet}, \code{CommandeService.listerCommandes} \\
  BF21 & Changer le statut d'une commande & P1 & Réalisé & \code{CommandeService.changerStatut} \\
  BF22 & Utilisateurs, activer ou désactiver & P2 & Réalisé & \code{AdminUtilisateurServlet}, \code{AuthService.changerEtatCompte} \\
  BF23 & Tableau de bord & P3 & Réalisé & \code{TableauDeBordServlet} \\
\end{xltabular}
}

\section{Règles de gestion}

{\small
\begin{xltabular}{\linewidth}{P{1.2cm}L}
  \caption{Mise en œuvre des règles de gestion}\label{tab:rg}\\
  \ligneentete
  \entete{Code} & \entete{Mise en œuvre} \\
  \endfirsthead
  \ligneentete
  \entete{Code} & \entete{Mise en œuvre} \\
  \endhead
  \bottomrule
  \endfoot
  RG01 & \code{AuthService} : recherche par e-mail avant l'inscription et la modification du profil ; contrainte d'unicité en base. \\
  RG02 & \code{AuthService.inscrire} force le rôle CLIENT ; le premier ADMIN est créé par \code{docs/sql/admin.sql}. \\
  RG03 & \code{CatalogueService.enregistrerProduit} exige une catégorie existante ; \code{@ManyToOne(optional = false)}. \\
  RG04 & \code{CatalogueService.supprimerCategorie} refuse si \code{CategorieDao.compterProduits} > 0 (produits inactifs compris). \\
  RG05 & \code{CatalogueService} : prix > 0 après arrondi au centime, stock ≥ 0 ; colonnes \code{DECIMAL(10,2)}. \\
  RG06 & Produit à stock 0 affiché avec le badge « Rupture » ; refusé par \code{PanierService}. \\
  RG07 & \code{PanierService} : quantité au panier + quantité demandée ≤ stock. \\
  RG08 & \code{AuthFilter} sur \code{/client/*} (panier et commandes) : rôle CLIENT obligatoire. \\
  RG09 & \code{CommandeService.passerCommande} refuse un panier vide. \\
  RG10 & Le prix du produit relu en base est copié dans \code{LigneCommande.prixUnitaire}. \\
  RG11 & Total calculé par le service à partir des lignes. \\
  RG12 & Décrément du stock et enregistrement dans une transaction ; rollback au premier stock insuffisant. \\
  RG13 & \code{CommandeService.statutsSuivants} : seules les transitions du cycle sont acceptées et proposées. \\
  RG14 & Réintégration des quantités à l'annulation (\code{ProduitDao.reintegrerStock}), même transaction. \\
  RG15 & \code{CatalogueService.supprimerProduit} désactive un produit déjà commandé ; le catalogue filtre \code{actif = true}. \\
  RG16 & \code{CommandeService.trouverCommandeClient} et \code{annulerCommandeClient} vérifient le propriétaire. \\
\end{xltabular}
}

\section{Besoins non fonctionnels}

{\small
\begin{xltabular}{\linewidth}{P{1.3cm}L}
  \caption{Mise en œuvre des besoins non fonctionnels}\label{tab:bnf}\\
  \ligneentete
  \entete{Code} & \entete{Mise en œuvre} \\
  \endfirsthead
  \ligneentete
  \entete{Code} & \entete{Mise en œuvre} \\
  \endhead
  \bottomrule
  \endfoot
  BNF01 & Mots de passe hachés avec jBCrypt (\code{BCrypt.hashpw}, \code{checkpw}). \\
  BNF02 & \code{AuthFilter} (\code{/client/*}, rôle CLIENT) et \code{AdminFilter} (\code{/admin/*}, rôle ADMIN). \\
  BNF03 & HQL à paramètres nommés uniquement ; affichage des données par \code{<c:out>}. \\
  BNF04 & Session invalidée à la déconnexion ; \code{session-timeout} de 30 minutes dans \code{web.xml}. \\
  BNF05 & Passage de commande et changement de statut dans une seule transaction Hibernate. \\
  BNF06 & Couches contrôleur, service, DAO, modèle ; aucun scriptlet dans les JSP. \\
  BNF07 & Messages de confirmation et d'erreur en français, contrôles HTML5 alignés sur le serveur, pages d'erreur 400, 403, 404, 405, 500. \\
  BNF08 & Bootstrap 5 : grille, tableaux \code{table-responsive}, barre de navigation repliable. \\
  BNF09 & Base \code{utf8mb4}, \code{EncodingFilter}, JSP en UTF-8 (\code{web.xml}), connexion JDBC en UTF-8. \\
  BNF10 & Pagination par 10 : catalogue, produits, commandes (client et administrateur), utilisateurs. \\
  BNF11 & Archive \code{boutique.war} déployée sur Tomcat 10.1 (Jakarta EE 10, espace de noms \code{jakarta.*}). \\
  BNF12 & Dépôt GitHub, une branche par fonctionnalité, fusion par pull request (tableau~\ref{tab:iterations}). \\
\end{xltabular}
}
